%%%%%%%%%%%%Outline%%%%%%%%%%%%%%%%%%

%%%
% message<a rejection returns one message that names the operation but under-determines the repair>
% eBPF runs networking, security, and observability in the kernel
% every program must clear a static verifier that rejects what it cannot prove safe
% a rejection returns one message that names the failing operation
% however the message under-determines the repair
% the developer fixes the rejection by trial and error

%%%
% message<we assemble bpfix-bench and find the message leaves the bug, the layer, and the fix site to recover>
% to learn what the message leaves out, we assemble bpfix-bench and find it leaves the bug, the layer, the fix site
% one message maps to many bugs: the most frequent covers 28 across nine mechanisms
% the message has no field for the layer, which can be source, compiler, environment, or verifier
% the flagged operation is the fix site in only 25.5%
% tools reading only the message inherit the gap
% such tools restate the message or regenerate the program

%%%
% message<the discarded reasoning survives in the log, so reconstructing the proof and locating the loss is principled>
% although the message omits all three, the verifier's reasoning survives in the log
% a rejection is a safety proof built instruction by instruction but not completed
% at the failing operation one needed fact cannot be proved
% the verifier then discards the proof and prints the symptom
% at log_level=2 the verifier records its abstract state after every instruction
% a tool reads that record to reconstruct the proof and find where a needed fact was lost
% reading the verifier's own states makes localization principled

%%%
% message<to recover the three we build \sys, a deterministic log-only tool, evaluated on localization and downstream repair>
% to recover the bug, the layer, and the fix site, we build \sys, a deterministic tool that reconstructs the proof and locates the fix site
% \sys reports the responsible layer, the weakest step, as best-effort
% we evaluate \sys two ways
% first, on bpfix-bench the labeled fix site lies within 13 instructions of the rejected line
% second, we feed the diagnostic to a downstream repair model, a separate LLM, and test repair lift
% on bpfix-test Qwen3.6 27B repairs 22 of 75 from raw and 38 from the diagnostic, 30 vs 44 with retry
% we read this repair lift as controlled usefulness on a suite still being extended

%%%
% message<three contributions: the finding, the tool, the downstream-repair result>
% we study rejections and show the message leaves the bug, the layer, and the fix site
% we build \sys, a log-only tool that recovers the discarded proof and locates the fix site
% we demonstrate \sys's diagnostic improves downstream LLM repair

%%%%%%%%%%%%Outline%%%%%%%%%%%%%%%%%%

\section{Introduction}
\label{sec:introduction}

eBPF now runs production networking, security, and observability inside the Linux kernel~\cite{gbadamosi2024ebpf}. Every program must first clear a static safety verifier, which rejects any program it cannot prove safe~\cite{gershuni2019simple}. A rejection returns one message that names the failing operation. However, the message under-determines the repair. The developer is left to fix the rejection by trial and error~\cite{deokar2024empirical}.

To learn what the message leaves out, we assemble \bench, a corpus of 235 reproduced rejections, and find that the message leaves the developer to recover the bug, the responsible layer, and the fix site. One message can map to many bugs; the most frequent covers 28 rejections across nine source mechanisms. The message carries no field for the responsible layer, which can be the developer's source, the compiler, the environment, or the verifier. The flagged operation is the fix site in only 25.5\% of cases. Tools that read only this message inherit the gap. Such tools either restate the message~\cite{rizza2025design} or regenerate the program to sidestep the rejection~\cite{zheng2024kgent,jia2025rex}.

Although the message omits all three, the verifier's reasoning survives in the log. A rejection is a safety proof the verifier builds instruction by instruction but cannot complete. At the failing operation one needed fact cannot be proved. The verifier then discards the proof and prints the symptom. Under \texttt{log\_level=2}, the verifier records its abstract state after every instruction~\cite{linuxkernel-bpf-verifier,vishwanathan2022sound}. A tool can read that record to reconstruct the discarded proof and find where a needed safety fact was lost. Reading the verifier's own states makes this localization principled.

To recover the bug, the layer, and the fix site that the message omits, we build \sys, a deterministic userspace tool that reconstructs the discarded proof from the log and locates the fix site. \sys also reports the responsible layer, the weakest step, as a best-effort heuristic. We evaluate \sys two ways. First, on \bench the labeled fix site lies within 13 instructions of the rejected line. Second, we feed the \sys diagnostic to a downstream repair model, a separate large language model, and test whether the diagnostic raises repair success. On \testset, a 75-case repair suite, Qwen3.6 27B repairs 22 of 75 from the raw log and 38 from the diagnostic in one shot, and 30 versus 44 with one retry. We read this repair lift as a controlled usefulness result on a suite still being extended.

\noindent\textbf{Contributions.} This paper makes the following contributions:
\begin{itemize}
\item We study eBPF verifier rejections and show that the message leaves the bug, the responsible layer, and the fix site for the developer to recover.
\item We build \sys, a log-only tool that recovers the verifier's discarded proof and locates the fix site.
\item We demonstrate that \sys's diagnostic improves downstream LLM repair.
\end{itemize}